\documentclass[10pt,a4paper]{amsart}
\usepackage{amsmath,amssymb,amsthm,graphicx}
% Retained preamble for the Stacks Project chapter snapshot in this directory.
% The chapter's own preamble.tex is a copy of the shared Stacks Project preamble; the block
% below is sliced MECHANICALLY, line for line, from it. Nothing was added, removed or
% redefined; the wrapper class, the xr-hyper/hyperref wiring and the externaldocument list
% are omitted because the assembler supplies the document class.
%
% lines taken from preamble.tex: 3, 5, 9, 11, 12, 156, 157, 158, ... (51 lines total)
%
% package coverage restored from the chapter's own preamble (2026-10-13): xy, xy-2cell

\usepackage[all]{xy}
\xyoption{2cell}
\UseAllTwocells

\newenvironment{reference}{\comment}{\endcomment}
\newenvironment{slogan}{\comment}{\endcomment}
\newenvironment{history}{\comment}{\endcomment}
\theoremstyle{plain}
\newtheorem{theorem}[subsection]{Theorem}
\newtheorem{proposition}[subsection]{Proposition}
\newtheorem{lemma}[subsection]{Lemma}
\theoremstyle{definition}
\newtheorem{definition}[subsection]{Definition}
\newtheorem{example}[subsection]{Example}
\newtheorem{exercise}[subsection]{Exercise}
\newtheorem{situation}[subsection]{Situation}
\theoremstyle{remark}
\newtheorem{remark}[subsection]{Remark}
\newtheorem{remarks}[subsection]{Remarks}
\numberwithin{equation}{subsection}
\def\lim{\mathop{\mathrm{lim}}\nolimits}
\def\colim{\mathop{\mathrm{colim}}\nolimits}
\def\Spec{\mathop{\mathrm{Spec}}}
\def\Hom{\mathop{\mathrm{Hom}}\nolimits}
\def\Ext{\mathop{\mathrm{Ext}}\nolimits}
\def\SheafHom{\mathop{\mathcal{H}\!\mathit{om}}\nolimits}
\def\SheafExt{\mathop{\mathcal{E}\!\mathit{xt}}\nolimits}
\def\Sch{\mathit{Sch}}
\def\Mor{\mathop{\mathrm{Mor}}\nolimits}
\def\Ob{\mathop{\mathrm{Ob}}\nolimits}
\def\Sh{\mathop{\mathit{Sh}}\nolimits}
\def\NL{\mathop{N\!L}\nolimits}
\def\CH{\mathop{\mathrm{CH}}\nolimits}
\def\proetale{{pro\text{-}\acute{e}tale}}
\def\etale{{\acute{e}tale}}
\def\QCoh{\mathit{QCoh}}
\def\Ker{\mathop{\mathrm{Ker}}}
\def\Im{\mathop{\mathrm{Im}}}
\def\Coker{\mathop{\mathrm{Coker}}}
\def\Coim{\mathop{\mathrm{Coim}}}
\DeclareMathSymbol{\boxtimes}{\mathbin}{AMSa}{"02}
\def\QCohstack{\mathcal{QC}\!\mathit{oh}}
\def\Cohstack{\mathcal{C}\!\mathit{oh}}
\def\Spacesstack{\mathcal{S}\!\mathit{paces}}
\def\Quotfunctor{\mathrm{Quot}}
\def\Hilbfunctor{\mathrm{Hilb}}
\def\Curvesstack{\mathcal{C}\!\mathit{urves}}
\def\Polarizedstack{\mathcal{P}\!\mathit{olarized}}
\def\Complexesstack{\mathcal{C}\!\mathit{omplexes}}
\def\Pic{\mathop{\mathrm{Pic}}\nolimits}
\def\Picardstack{\mathcal{P}\!\mathit{ic}}
\def\Picardfunctor{\mathrm{Pic}}
\def\Deformationcategory{\mathcal{D}\!\mathit{ef}}
\usepackage{cleveref}
\crefname{originalproposition}{original proposition}{original propositions}
\Crefname{originalproposition}{Original proposition}{Original propositions}
\makeatletter
\newlabel{derived-lemma-fully-faithful-adjoint-kernel-zero}{{09J1}{source}{}{sourcepointer-1}{}}
\newlabel{derived-lemma-fully-faithful-adjoint-kernel-zero@cref}{{[originalproposition][0][]09J1}{[1][0][]source}{}{}{}}
\makeatother
\crefname{originalproposition}{original proposition}{original propositions}
\Crefname{originalproposition}{Original proposition}{Original propositions}
\makeatletter
\newlabel{dpa-remark-pro-system-koszul}{{0GZ5}{source}{}{sourcepointer-2}{}}
\newlabel{dpa-remark-pro-system-koszul@cref}{{[originalproposition][0][]0GZ5}{[1][0][]source}{}{}{}}
\makeatother
\crefname{originalproposition}{original proposition}{original propositions}
\Crefname{originalproposition}{Original proposition}{Original propositions}
\makeatletter
\newlabel{lemma-cartesian-qis-equivalence}{{0GZF}{source}{}{sourcepointer-3}{}}
\newlabel{lemma-cartesian-qis-equivalence@cref}{{[originalproposition][0][]0GZF}{[1][0][]source}{}{}{}}
\makeatother
\crefname{originalproposition}{original proposition}{original propositions}
\Crefname{originalproposition}{Original proposition}{Original propositions}
\makeatletter
\newlabel{lemma-different-tensors}{{0GZJ}{source}{}{sourcepointer-4}{}}
\newlabel{lemma-different-tensors@cref}{{[originalproposition][0][]0GZJ}{[1][0][]source}{}{}{}}
\makeatother
\crefname{originalproposition}{original proposition}{original propositions}
\Crefname{originalproposition}{Original proposition}{Original propositions}
\makeatletter
\newlabel{lemma-pro-isomorphic}{{0GZH}{source}{}{sourcepointer-5}{}}
\newlabel{lemma-pro-isomorphic@cref}{{[originalproposition][0][]0GZH}{[1][0][]source}{}{}{}}
\makeatother
\crefname{originalproposition}{original proposition}{original propositions}
\Crefname{originalproposition}{Original proposition}{Original propositions}
\makeatletter
\newlabel{more-algebra-definition-derived-complete}{{091S}{source}{}{sourcepointer-6}{}}
\newlabel{more-algebra-definition-derived-complete@cref}{{[originalproposition][0][]091S}{[1][0][]source}{}{}{}}
\makeatother
\crefname{originalproposition}{original proposition}{original propositions}
\Crefname{originalproposition}{Original proposition}{Original propositions}
\makeatletter
\newlabel{more-algebra-lemma-naive-derived-completion}{{091W}{source}{}{sourcepointer-7}{}}
\newlabel{more-algebra-lemma-naive-derived-completion@cref}{{[originalproposition][0][]091W}{[1][0][]source}{}{}{}}
\makeatother
\crefname{originalproposition}{original proposition}{original propositions}
\Crefname{originalproposition}{Original proposition}{Original propositions}
\makeatletter
\newlabel{more-algebra-proposition-noetherian-naive-completion-is-completion}{{0922}{source}{}{sourcepointer-8}{}}
\newlabel{more-algebra-proposition-noetherian-naive-completion-is-completion@cref}{{[originalproposition][0][]0922}{[1][0][]source}{}{}{}}
\makeatother
\title{Stacks, Differential Graded Sheaves: proposition-derived-complete}
\author{The Stacks Project Authors}
\date{stacks-project:sdga.tex:local-snapshot-20261012; GNU FDL 1.2 or later}
\begin{document}
\maketitle
\noindent\textit{Editorial note.} Statement span is lines 6601-6618 of sources/stacks-sdga/pinned-original-source.tex; proof span is lines 6620-6722, the complete author proof (101 body lines).\par\medskip
\noindent\textit{Editorial note.} Navigation stubs: 8 cited-but-not-included labels are supplied as non-load-bearing pointers with their pinned tags where the snapshot tag index has them.\par\medskip
\noindent\textit{Editorial note.} extra\_packages is empty on purpose: the retained preamble already loads xy with options and a second declaration clashes in the assembler wrapper.\par\medskip
\noindent\textit{Editorial note.} Text is retained verbatim, including the author's own typos and placeholders.\par\medskip
\noindent\textit{Editorial note.} Queue index 2 of queue\_sdga.json (pb=4129, stmt 6601-6618, proof 6620-6722).\par\medskip
\noindent\textit{Allowed standard prerequisites (editorial).}\begin{itemize}
\item Explicit prerequisite boundary. The cited inputs are prerequisites; the argument in the retained span is not replaced by any cited result.
\end{itemize}\medskip
\ifdefined\hypertarget\hypertarget{sourcepointer-1}{}\fi
\noindent\textit{Source reference.} Original proposition 09J1: sources/stacks-sdga/pinned-original-source.tex, source lines 6722--6722. Cited by the retained statement or proof but not part of the collected spans; supplied only as a navigation pointer. The argument does not depend on its content.\par\medskip
\ifdefined\hypertarget\hypertarget{sourcepointer-2}{}\fi
\noindent\textit{Source reference.} Original proposition 0GZ5: sources/stacks-sdga/pinned-original-source.tex, source lines 6722--6722. Cited by the retained statement or proof but not part of the collected spans; supplied only as a navigation pointer. The argument does not depend on its content.\par\medskip
\ifdefined\hypertarget\hypertarget{sourcepointer-3}{}\fi
\noindent\textit{Source reference.} Original proposition 0GZF: sources/stacks-sdga/pinned-original-source.tex, source lines 6358--6368. Cited by the retained statement or proof but not part of the collected spans; supplied only as a navigation pointer. The argument does not depend on its content.\par\medskip
\ifdefined\hypertarget\hypertarget{sourcepointer-4}{}\fi
\noindent\textit{Source reference.} Original proposition 0GZJ: sources/stacks-sdga/pinned-original-source.tex, source lines 6557--6560. Cited by the retained statement or proof but not part of the collected spans; supplied only as a navigation pointer. The argument does not depend on its content.\par\medskip
\ifdefined\hypertarget\hypertarget{sourcepointer-5}{}\fi
\noindent\textit{Source reference.} Original proposition 0GZH: sources/stacks-sdga/pinned-original-source.tex, source lines 6452--6459. Cited by the retained statement or proof but not part of the collected spans; supplied only as a navigation pointer. The argument does not depend on its content.\par\medskip
\ifdefined\hypertarget\hypertarget{sourcepointer-6}{}\fi
\noindent\textit{Source reference.} Original proposition 091S: sources/stacks-sdga/pinned-original-source.tex, source lines 6722--6722. Cited by the retained statement or proof but not part of the collected spans; supplied only as a navigation pointer. The argument does not depend on its content.\par\medskip
\ifdefined\hypertarget\hypertarget{sourcepointer-7}{}\fi
\noindent\textit{Source reference.} Original proposition 091W: sources/stacks-sdga/pinned-original-source.tex, source lines 6722--6722. Cited by the retained statement or proof but not part of the collected spans; supplied only as a navigation pointer. The argument does not depend on its content.\par\medskip
\ifdefined\hypertarget\hypertarget{sourcepointer-8}{}\fi
\noindent\textit{Source reference.} Original proposition 0922: sources/stacks-sdga/pinned-original-source.tex, source lines 6722--6722. Cited by the retained statement or proof but not part of the collected spans; supplied only as a navigation pointer. The argument does not depend on its content.\par\medskip
\begin{proposition}
\label{proposition-derived-complete}
Let $R$ be a Noetherian ring. Let $I \subset R$ be an ideal.
The following three categories are canonically equivalent:
\begin{enumerate}
\item Let $\mathcal{A}$ be the sheaf of $R$-algebras on $\mathbf{N}$
corresponding to the inverse system of $R$-algebras $A_n = R/I^n$.
The category $\mathit{QC}(\mathcal{A})$.
\item Choose generators $f_1, \ldots, f_r$ of $I$. Let
$\mathcal{B}$ be the sheaf of differential graded $R$-algebras
on $\mathbf{N}$ corresponding to the inverse system of
Koszul algebras on $f_1^n, \ldots, f_r^n$.
The category $\mathit{QC}(\mathcal{B})$.
\item The full subcategory $D_{comp}(R, I) \subset D(R)$
of derived complete objects, see More on Algebra, Definition
\ref{more-algebra-definition-derived-complete} and text following.
\end{enumerate}
\end{proposition}

\begin{proof}
Consider the obvious morphism
$f : (\Sh(\mathbf{N}), \mathcal{A}) \to (\Sh(pt), R)$
of ringed topoi
and let us consider the adjoint functors $Lf^*$ and $Rf_*$.
The first restricts to a functor
$$
F : D_{comp}(R, I) \longrightarrow \mathit{QC}(\mathcal{A})
$$
which sends an object $K$ of $D_{comp}(R, I)$ represented by a
K-flat complex $K^\bullet$ to the object $(K^\bullet \otimes_R R/I^n)$
of $\mathit{QC}(\mathcal{A})$. The second restricts to a functor
$$
G : \mathit{QC}(\mathcal{A}) \longrightarrow D_{comp}(R, I)
$$
which sends an object $(M_n^\bullet)$ of $\mathit{QC}(\mathcal{A})$
to $R\lim M_n^\bullet$. The output is derived complete for example
by More on Algebra, Lemma \ref{more-algebra-lemma-naive-derived-completion}.
Also, it follows from More on Algebra, Proposition
\ref{more-algebra-proposition-noetherian-naive-completion-is-completion}
that $G \circ F = \text{id}$. Thus to see that $F$ and $G$
are quasi-inverse equivalences it suffices to see that the kernel
of $G$ is zero (see Derived Categories, Lemma
\ref{derived-lemma-fully-faithful-adjoint-kernel-zero}).
However, it does not appear easy to show this directly!

\medskip\noindent
In this paragraph we will show that $\mathit{QC}(\mathcal{A})$ and
$\mathit{QC}(\mathcal{B})$ are equivalent.
Write $\mathcal{B} = (B_n)$ where $B_n$ is the Koszul complex
viewed as a cochain complex in degrees $-r, -r + 1, \ldots, 0$.
By Divided Power Algebra, Remark \ref{dpa-remark-pro-system-koszul}
(but with chain complexes
turned into cochain complexes)
we can find $1 < n_1 < n_2 < \ldots$
and maps of differential graded $R$-algebras
$B_{n_i} \to E_i \to R/(f_1^{n_i}, \ldots, f_r^{n_i})$
and $E_i \to B_{n_{i - 1}}$ such that
$$
\xymatrix{
B_{n_1} \ar[d] &
B_{n_2} \ar[d] \ar[l] &
B_{n_3} \ar[d] \ar[l] &
\ldots \ar[l] \\
E_1 \ar[d] &
E_2 \ar[l] \ar[d] &
E_3 \ar[l] \ar[d] &
\ldots \ar[l] \\
B_1 &
B_{n_1} \ar[l] &
B_{n_2} \ar[l] &
\ldots \ar[l]
}
$$
is a commutative diagram of differential graded $R$-algebras
and such that $E_i \to R/(f_1^{n_i}, \ldots, f_r^{n_i})$ is
a quasi-isomorphism. We conclude
\begin{enumerate}
\item there is an equivalence between $\mathit{QC}(\mathcal{B})$ and
$\mathit{QC}((E_i))$,
\item there is an equivalence between
$\mathit{QC}((E_i))$ and $\mathit{QC}((R/(f_1^{n_i}, \ldots, f_r^{n_i})))$,
\item there is an equivalence between
$\mathit{QC}((R/(f_1^{n_i}, \ldots, f_r^{n_i})))$ and
$\mathit{QC}(\mathcal{A})$.
\end{enumerate}
Namely, for (1) we can apply Lemma \ref{lemma-pro-isomorphic}
to the diagram above which shows that $(E_i)$ and $(B_n)$ are pro-isomorphic.
For (2) we can apply  Lemma \ref{lemma-cartesian-qis-equivalence} to the
inverse system of quasi-isomorphisms $E_i \to R/(f_1^{n_i}, \ldots, f_r^{n_i})$.
For (3) we can apply Lemma \ref{lemma-pro-isomorphic}
and the elementary fact that the inverse systems $(R/I^n)$ and
$(R/(f_1^{n_i}, \ldots, f_r^{n_i})$ are pro-isomorphic.

\medskip\noindent
Exactly as in the first paragraph of the proof we can define
adjoint functors\footnote{It can be shown that these functors are,
via the equivalences above, compatible with $F$ and $G$ defined before.}
$$
F' : D_{comp}(R, I) \longrightarrow \mathit{QC}(\mathcal{B})
\quad\text{and}\quad
G' : \mathit{QC}(\mathcal{B}) \longrightarrow D_{comp}(R, I).
$$
The first sends an object $K$ of $D_{comp}(R, I)$ represented by a
K-flat complex $K^\bullet$ to the object $(K^\bullet \otimes_R B_n)$
of $\mathit{QC}(\mathcal{B})$. The second sends an object
$(M_n)$ of $\mathit{QC}(\mathcal{B})$ to $R\lim M_n$.
Arguing as above it suffices to show that the kernel of $G'$ is zero.
So let $\mathcal{M} = (M_n)$ be a good sheaf of differential graded
modules over $\mathcal{B}$ which represents an object of
$\mathit{QC}(\mathcal{B})$ in the kernel of $G'$. Then
$$
0 = R\lim M_n \Rightarrow
0 = (R\lim M_n) \otimes_R^\mathbf{L} B_t =
R\lim (M_n \otimes_R^\mathbf{L} B_t)
$$
By Lemma \ref{lemma-different-tensors} we have
$R\lim (M_n \otimes_R^\mathbf{L} B_t) =
R\lim (M_n \otimes_{B_n}^\mathbf{L} B_t)$.
Since $(M_n)$ is an object of $\mathit{QC}(\mathcal{B})$ we
see that the inverse system $M_n \otimes_{B_n}^\mathbf{L} B_t$
is eventually constant with value $M_t$. Hence $M_t = 0$ as desired.
\end{proof}

\end{document}
